\subsection{Modelos de aprendizado de máquina}
\label{subsec:modelos}

Foram avaliados sete classificadores individuais e um ensemble para examinar
se a discriminação entre os grupos dependia de uma família de modelos. Extra
Trees (ET) e Random Forest (RF) representam conjuntos de árvores com
fronteiras não lineares; a rede convolucional temporal (TCN) usa a sequência
das leituras; a rede perceptron multicamadas (MLP) modela relações não
lineares sem convoluções temporais; o k-NN classifica por proximidade entre
leituras; e SVM linear e análise discriminante linear (LDA) servem como
referências lineares. Assim, os modelos testam diferentes formas de usar as
mesmas entradas, sem pressupor que uma arquitetura seja a mais adequada.

A comparação LOSO usou 55.946 leituras de 31 coletas pertencentes a dez
vasos, após as exclusões descritas anteriormente. Foram considerados dois
conjuntos fixos de entradas: (i) MQ2, MQ3, MQ8, MQ135 e MQ138; e (ii) esses
cinco sinais acrescidos do índice de umidade do solo, da temperatura e da
pressão. A classe 0 corresponde aos vasos inoculados, e a classe 1, aos
vasos controle.

As configurações da Tabela~\ref{tab:configuracao_modelos} foram fixadas para
esta comparação exploratória. Em uma etapa anterior do projeto, quatro
configurações candidatas de ET foram comparadas no conjunto de treinamento
por validação cruzada agrupada por coleta. Entre elas, a configuração com
\texttt{max\_features="sqrt"}, mínimo de dez leituras por folha e sem pesos
de classe apresentou a maior acurácia balanceada média. Essa configuração
foi mantida; as 900 árvores vieram do protocolo anterior e não foram
selecionadas por aquela busca. Os demais modelos receberam configurações
únicas registradas na implementação. Não houve busca sistemática de
hiperparâmetros para os sete modelos sob o protocolo LOSO; portanto, a
comparação examina famílias de classificadores com configurações fixas,
sem estimar a melhor parametrização possível de cada família.

Na validação \textit{leave-one-subject-out} (LOSO), cada modelo foi ajustado
novamente em cada uma das dez rodadas, reservando-se um vaso inteiro para
avaliação. Todas as leituras de um vaso permaneceram no mesmo lado da divisão.
ET e RF receberam as entradas diretamente. Para k-NN, MLP, SVM e LDA, a
padronização foi ajustada somente nas leituras de treinamento da rodada;
a TCN também usou padronização ajustada apenas no treinamento. A calibração
probabilística da SVM empregou três divisões internas agrupadas por vaso,
restritas aos nove vasos de treinamento da rodada externa. Não houve
subamostragem ou geração de leituras sintéticas para balancear as classes;
os tratamentos de desbalanceamento usados por RF, TCN e SVM estão indicados
na Tabela~\ref{tab:configuracao_modelos}.

\begin{table}[H]
\centering
\footnotesize
\caption{Configurations of the evaluated classifiers.}
\label{tab:configuracao_modelos}
\begin{tabularx}{\textwidth}{@{}p{1.35cm}X@{}}
\toprule
\textbf{Model} & \textbf{Configuration} \\
\midrule
ET & 900 trees; \texttt{max\_features="sqrt"}; at least 10 readings per
leaf; \texttt{bootstrap=False}; random seed 42. \\
RF & 300 trees; \texttt{class\_weight="balanced"}; bootstrap sampling
(library default); random seed 42. \\
TCN & Three temporal residual blocks with 24, 24, and 32 channels; two
convolutions per block, kernel size 3, dilations 1, 2, and 4, global
average pooling, and dropout rate 0.20. Windows of 128
readings with stride 64; 60 epochs, batch size 32 windows, AdamW
($10^{-3}$ learning rate; $10^{-4}$ weight decay), and class-weighted
cross-entropy; random seed 42. \\
MLP & Two hidden layers with 64 and 32 neurons and ReLU activation;
\texttt{alpha=0.001}; batch size 2048; up to 250 epochs, with early stopping
using 15\% of the training readings and a patience of 12 epochs;
random seed 42. \\
k-NN & Five neighbors; Euclidean distance and distance-weighted voting. \\
SVM & Linear SVM with $C=1$, balanced class weights, and up to 10,000
iterations; sigmoid probability calibration using three internal splits
grouped by pot. \\
LDA & Linear discriminant analysis with \texttt{solver="lsqr"} and
automatic covariance shrinkage. \\
Ensemble & Unweighted mean of the class-1 probabilities produced by ET, RF,
TCN, MLP, and k-NN; classification threshold of 0.5. \\
\bottomrule
\end{tabularx}
\end{table}

No ET, o sorteio de um subconjunto de atributos em cada divisão
(\texttt{max\_features="sqrt"}) aumenta a diversidade entre as árvores,
enquanto o mínimo de dez leituras por folha restringe divisões baseadas
em poucos registros. O ET usa todas as leituras de treinamento em cada
árvore (\texttt{bootstrap=False}); o RF usa reamostragem com reposição,
que constitui uma diferença metodológica entre os dois conjuntos de
árvores. No RF, \texttt{class\_weight="balanced"} ajusta a contribuição
das classes na construção das árvores segundo suas frequências no treino.
As 300 árvores do RF foram mantidas como configuração fixa da comparação,
sem estudo de sensibilidade para esse número.

Na TCN, as dilatações 1, 2 e 4 ampliam o campo receptivo sem aumentar
proporcionalmente o número de parâmetros. Com duas convoluções de núcleo
3 por bloco, o campo receptivo local é de 29 leituras; a agregação média
global reúne as respostas ao longo da janela completa. O \emph{dropout}
de 0,20 e o decaimento de pesos do AdamW são mecanismos de regularização.
A janela de 128 leituras fornece contexto temporal, e o passo de 64 cria
sobreposição nominal de 50\%; esses tamanhos, assim como as 60 épocas,
foram mantidos da implementação anterior e não foram otimizados na
comparação LOSO.

A MLP de duas camadas ocultas constitui uma referência neural sem
convoluções. A penalização $L_2$ (\texttt{alpha=0.001}) e a parada
antecipada limitam o ajuste aos dados de treinamento; o lote de 2048
leituras reduz o número de atualizações por época nesta base. No k-NN,
cinco vizinhos fornecem uma decisão local, com maior influência dos
registros mais próximos pela ponderação inversa da distância. Para a SVM,
$C=1$ fixa o compromisso entre margem e erros de classificação, enquanto
os pesos equilibrados consideram a diferença de frequência entre classes;
a calibração sigmoide fornece probabilidades para comparação no mesmo
limiar dos demais modelos. Na LDA, a regularização automática da matriz
de covariância é útil diante de entradas potencialmente correlacionadas.
Os valores numéricos dessas configurações foram fixados como referências
de comparação e não constituem resultados de uma otimização conjunta.

Para a TCN, as leituras foram ordenadas por coleta e tempo e organizadas em
janelas de 128 leituras, com passo de 64, sem atravessar os limites entre
coletas. Uma janela adicional foi alinhada ao fim da coleta quando o passo
regular não cobria suas últimas leituras. Se uma coleta tivesse menos de
128 leituras, a última leitura seria repetida até completar a janela.
Cada janela teve cinco ou oito canais, conforme o conjunto de entradas, e
recebeu a classe da respectiva coleta. Os pesos da função de perda foram
calculados a partir da frequência das classes nas janelas de treinamento
de cada rodada. A rede produziu uma probabilidade da classe 1 por janela;
para cada leitura, calculou-se a média das probabilidades das janelas que
a continham. As 60 épocas da TCN foram fixas, sem parada antecipada; a
MLP, por sua vez, usou parada antecipada conforme a
Tabela~\ref{tab:configuracao_modelos}. A implementação foi denominada TCN,
pois não incluiu mecanismo de atenção nem explicabilidade integrada.

O ensemble foi calculado depois das predições dos modelos individuais, sem
treinar um classificador adicional. Para cada leitura do vaso reservado,
sejam $p_{\mathrm{ET}}$, $p_{\mathrm{RF}}$, $p_{\mathrm{TCN}}$,
$p_{\mathrm{MLP}}$ e $p_{\mathrm{kNN}}$ as probabilidades atribuídas à classe 1
(controle).
A decisão combinada foi
\begin{equation}
    \widehat{p}_1=\frac{p_{\mathrm{ET}}+p_{\mathrm{RF}}+
    p_{\mathrm{TCN}}+p_{\mathrm{MLP}}+p_{\mathrm{kNN}}}{5},
    \qquad \widehat{y}=\mathbb{1}(\widehat{p}_1\geq 0{,}5).
\end{equation}
Portanto, trata-se de média de probabilidades com pesos iguais, não de
votação majoritária das classes previstas. Os cinco componentes reúnem
métodos arbóreos, temporal, neural e de vizinhança. SVM e LDA foram mantidos
como referências individuais; seus resultados inferiores nessa base
motivaram sua ausência na combinação final. A escolha de cinco componentes
privilegiou a variedade de famílias, sem selecionar apenas os modelos de
maior acurácia. Não foi feita uma comparação sistemática com combinações de
dois ou três classificadores; portanto, não se afirma que esta composição
seja superior a elas. A composição do ensemble, as exclusões de coletas e
a seleção das entradas foram definidas após inspeção de resultados nos
mesmos dez vasos. Por isso, esta comparação LOSO é exploratória e não
constitui uma estimativa independente de desempenho em novos vasos.
